% year: תשפ"ד
% semester: ב
% moed: א
% number: 4ב
% categories: taylor
% source: exams/מבחנים/תשפד/מועד א תשפד סמסטר ב.pdf

\begin{question}
מצאו קירוב לינארי של $\sqrt{15}$
\end{question}

\begin{hint}
השתמשו במשיק לגרף $g(x)=\sqrt x$ בנקודה קרובה ל-$15$ שבה השורש ידוע.
\end{hint}

\begin{solution}
נשתמש בפונקציה $g(x)=\sqrt x$ סביב הנקודה $a=16$ (הקרובה ל-$15$, שבה השורש ידוע). $g(16)=4$, $g'(x)=\frac{1}{2\sqrt x}$, $g'(16)=\frac18$. הקירוב הלינארי (משוואת המשיק):
\[ g(x)\approx g(16)+g'(16)(x-16)=4+\frac{x-16}{8}. \]
עבור $x=15$: $\sqrt{15}\approx4-\frac18=\boxed{3.875}$.
(לשם השוואה $\sqrt{15}\approx3.873$. מאחר ש-$g''<0$ — $g$ קעורה — המשיק נמצא מעל הגרף, ולכן זהו קירוב מלמעלה; לפי שארית לגרנז' השגיאה היא $\frac{|g''(c)|}{2}\cdot1^2=\frac{1}{8c^{3/2}}<\frac{1}{8\cdot 15^{3/2}}<0.003$.)
\end{solution}
